\section{Positive source coordinates and a congruence inside a Pell root}
\label{sec:positive-coordinate-carrier}

The new equation has thirteen strictly positive unknowns, exact total degree
$28$, and $98$ live binary arithmetic operations.  The degree-$34$,
$93$-operation equation of package~61 is the comparison point.  The change
is an existential coordinate transformation, not a same-coordinate identity.
The inherited native compiler and Pell-index theorems remain explicitly
identified dependencies.

The central choice is to supply the positive gap $C$ instead of the native
index $J$.  Set $q=F+\alpha+C$; the source bounds now hold as soon as the
fixed congruence on $q$ is established.  That congruence is placed inside
the normalized strong Pell root.  A monic square equation recovers the
integrality of that root, and therefore the congruence.  This permits the
transport quotient's sign to be recovered at the end, removing the
predecessor's degree-two positivity multiplier.

\subsection{Literal integer polynomial}

The positive variables are
\[
 F,\alpha,C,f_0,i,T,X,Y,r,\eta,\zeta,u,\gamma.
\]
Use the whole compact source-44 tuple
$(\beta,K_0,M_C,M_F)$, with $\beta=2^{590069}-1$ and the already shifted
high mask $M_F$.  Put
\[
 B=\beta+1,\qquad b=3\beta,\qquad L=\beta b^2.
\]
These are fixed numerals.  The following names are polynomial abbreviations,
not additional witnesses:
\begin{align*}
 q&=F+\alpha+C,& k&=\eta+\zeta,&c&=kY+\eta,\\
 E&=XY,&\ell&=EkY,&a&=E+Y,\\
 H&=4a+3,&\Delta&=a^2+H=(a+2)^2-1,&I&=ic^2,\\
 N_f&=bf_0+q-B,&F_2&=N_f^2,&M_0&=q^3(q-1),\\
 S_\beta&=4\beta qY,&M_\beta&=qS_\beta.&&
\end{align*}
The remaining registers are
\begin{align*}
 G&=(q-1)(q-F)+\alpha,\\
 \widehat R&=\beta G(q^2-1)+(M_C+qM_F)(q-1),\\
 B_t&=(qK_0+X)C-qF,\\
 Z&=M_0(\ell+r)+(q-1)Y-q^3B_t,\\
 \widehat D&=M_\beta(X+ac)
       +H\{M_\beta\gamma+S_\beta X-\beta(k-1)+\widehat R\},\\
 \widehat V&=F_2(\beta cT-\widehat R)-Lc,\\
 \widehat Y&=\widehat V+Lu.
\end{align*}
Define four integer residuals
\begin{align}
 A_0&=Z^2-M_0^2\{\ell(\ell+k)+1\},\label{eq:new-first}\\
 A_1&=\widehat D^{\,2}-M_\beta^2(1+\Delta c^2),\label{eq:new-main}\\
 A_a&=b^2\widehat Y^{\,2}
       -\Delta(F_2-b^2)(\widehat Y^{\,2}-\widehat V^{\,2})-b^2L^2,
       \label{eq:new-aux}\\
 A_s&=F_2-b^2(1+\Delta I^2).\label{eq:new-strong}
\end{align}
The equation is
\begin{equation}
 \boxed{\;P_{13}=A_0^2+A_1^2+A_a^2+A_s^2=0.\;}
 \label{eq:new-p13}
\end{equation}
Every coefficient in this definition is an integer.  There are no variable
exponents or unspecified program parameters.

For the proof only, use the rational abbreviations
\begin{align*}
 f&=N_f/b,&J&=(q-1)/\beta,&R&=\widehat R/\beta,\\
 S&=4qY,&M&=qS=4q^2Y,&v&=k-R-1,\\
 D_*&=M(X+ac)+H(M\gamma+SX-v),&&&\\
 V&=f^2(cT-R)-c,&U&=\Delta(f^2-1).&&
\end{align*}
Then $\widehat D=\beta D_*$, $\widehat V=LV$, and the residuals have the
exact interpretations
\begin{align*}
 A_s&=b^2(f^2-\Delta I^2-1),\\
 A_1&=\beta^2\{D_*^2-M^2(1+\Delta c^2)\},\\
 A_a&=b^2L^2\{(V+u)^2-Uu(2V+u)-1\}.
\end{align*}
These identities clear fixed denominators; they do not introduce divisions
into the delivered polynomial.

\subsection{Recovery of the native congruence}

Suppose \eqref{eq:new-p13} holds.  All four residuals vanish separately.
As $q\ge3$ and $f_0\ge1$,
\[
 N_f\ge2\beta+2>0.
\]
Equation \eqref{eq:new-strong} makes $(N_f/b)^2=1+\Delta I^2$ an integer.
A rational square root of an integer is integral.  Hence $f=N_f/b$ is a
positive integer and
\[
 q-\beta-1=3\beta j,\qquad J=3j+1
\]
for an integer $j$.  The inequality $q>0$ excludes $j<0$.  Thus $J\ge1$,
$q\ge B\ge16$, and, since $\beta\equiv1\pmod3$,
$q\equiv2\pmod3$.  Also $R$ is now integral with its exact native formula.

All of $F,\alpha,C$ are positive and $q=F+\alpha+C$.  The authentic mask
bounds give
\begin{equation}
 3q<R<q^4-2.\label{eq:new-rbounds}
\end{equation}
For example, $G=q\alpha+(q-1)C$ lies between $q$ and $q^2-q$, while
$0<(M_C+qM_F)J<(q-1)(1+2q)$.  These estimates use no transport quotient.

\subsection{The first carrier and normalized rank recovery}

Equation \eqref{eq:new-first} makes $\tau=Z/M_0$ an integer square root of
$\ell(\ell+k)+1$, possibly with negative sign.  Reducing $M_0\mid Z$
modulo $q^3$ and $q-1$ gives
\[
 q^3\mid Y,\qquad q-1\mid B_t.
\]
Write $s=Y/q^3>0$ and $t=B_t/(q-1)\in\mathbb Z$.  Thus
\[
 \tau=\ell+r+s-t,\qquad \tau^2-\ell(\ell+k)=1.
\]
Neither $t$ nor $\tau$ is assumed positive at this stage.

Similarly, \eqref{eq:new-main} gives an integer
$g=D_*/M$ with $g^2-\Delta c^2=1$.  It is positive: $v<k$, $c>kY$,
$H<5a$, and $MY>5$ imply
\[
 D_*>Mac-Hk>ak(MY-5)>0.
\]
Put $A=a+2$ and $P_0=2XY^2+1$.  Pell classification supplies $n,p\ge1$:
\[
 |\tau|=\chi_{P_0}(n),\quad k=2\psi_{P_0}(n),\qquad
 g=\chi_A(p),\quad c=\psi_A(p).
\]
The usual recurrence gives $P_0>A$; since $c>k$, monotonicity implies
$n<p$.  If $n=1$, then $k=2$, $\eta=\zeta=1$, and $c=2Y+1$, whereas
$p\ge2$ forces $c\ge2A\ge4Y+4$.  Hence $n\ge2$, and
\[
 c>kY\ge4P_0Y>8XY^3\ge8q^9>2R,\qquad c>2p.
\]

The \emph{normalized} strong equation is essential here.  It gives
$f=\chi_A(m)$ and $I=\psi_A(m)$.  Since $I=ic^2$, strong divisibility
gives $p\mid m$.  Write $m=pv_0$.  Composition and reduction modulo $c$
give
\[
 I/c=\psi_g(v_0)\equiv v_0g^{v_0-1}\pmod c.
\]
Because $g^2\equiv1\pmod c$, it follows that $c\mid v_0$, hence
$c\mid m$ and $m\ge pc>2p$.  This argument does not require the
large-$c$ hypothesis of the separate ordinary-strong rank lemma.

Set $M_a=\Delta I$.  Then $M_a^2=U$, $c\mid M_a$, and
\eqref{eq:new-aux} becomes
\[
 M_a^2V^2-(M_a^2-1)(V+u)^2=1.
\]
Both displayed Pell coordinates are nonzero.  The definition of $V$ and
$f^2\equiv1\pmod c$ give
\[
 V\equiv-c\pmod f,\qquad V\equiv-R\pmod c.
\]
All hypotheses of the inherited signed auxiliary-index theorem are now
established, including $0<p,R<c/2$ and $c\mid m$.  That theorem yields
$p=R$.  No deleted quotient or transport sign has been used.

\subsection{The adaptive congruence and transport sign}

The main carrier gives $M\mid H(SX-v)$.  Since $q\mid Y$ and
$q\equiv2\pmod3$, we have $\gcd(q,H)=1$.  Moreover, $H$ is odd and
$\gcd(Y,H)\mid3$, so
\[
 \mu=\frac{S}{\gcd(S,H)}\ge\frac{4qY}{3}>q^4>R+1.
\]
Reduction modulo $S$ gives $\mu\mid v$.  Both $S$ and $M$ divide
$4XY^2$ already, because $q^3\mid Y$; no assumption $q\mid X$ is needed.
The first Pell recurrence gives $k\equiv2n\pmod{4XY^2}$.  Thus
$\mu\mid2n-R-1$.  As $n<p=R$ and $|2n-R-1|<R<\mu$,
\[
 2n=R+1.
\]
Therefore $v=k-2n$ is a positive multiple of $4XY^2$, in particular of
$M$.  The main congruence now gives $q\mid HX$, hence $q\mid X$.
The integers
\[
 w=X/q>0,\qquad h=v/E>0,\qquad
 \gamma_{\rm old}=\gamma+w-v/M
\]
restore the source scale and index quotient.  The last one is positive:
$n\ge2$ implies $c>X$, while $g>(a+1)c$, so
$H\gamma_{\rm old}=g-ac-X>c-X>0$.

After $X=qw$, $B_t=q((K_0+w)C-F)$.  As $q-1\mid B_t$ and
$\gcd(q,q-1)=1$, the number
\[
 z_{\rm old}=\frac{(K_0+w)C+q-F-1}{q-1}
\]
is integral.  Its numerator is positive, since $K_0,w,C>0$ and
$q-F=\alpha+C\ge2$.  Thus $z_{\rm old}>0$ and the authentic transport
equation is restored.  This argument allows $z_{\rm old}-1=0$ and does
not assume $w>q$ for arbitrary solutions.

Finally $T_{\rm old}=fT$ and $y=|V+u|$ restore the remaining positive
source coordinates.  The inherited native compiler extracts a finite
syntactic ZFC refutation.

\subsection{Constructive completeness}

Conversely, start with canonical positive normalized source witnesses for
a refutation.  Supply their $F,\alpha,C$ and put
\[
 j=(J-1)/3,\qquad f_0=f-j.
\]
The first is integral, and $f_0>0$ because
$f\ge a+2>Y\ge q^3>q>j$.  This gives $N_f=bf$ and $A_s=0$.
The canonical odd-index reconstruction from package~61 supplies fresh
$T,u>0$ with
\[
 V_* =\frac{\chi_{M_a}(R)}{M_a},\quad y_* =\psi_{M_a}(R),\quad
 T=\frac{V_*+c+Rf^2}{cf^2},\quad u=y_*-V_*.
\]
Its integrality uses the established squared-modulus congruences and
$R\equiv3\pmod4$.  It gives $A_a=0$.

Let $t=q(z_{\rm old}-1)\ge0$ and set
\[
 r=\tau-\ell-s+t,\qquad
 \gamma=\gamma_{\rm old}+v/M-w.
\]
Both are positive.  For the first, $n\ge2$ gives
$k\ge8XY^2+4>3s$, and
\[
 \tau^2-(\ell+k/3)^2
 =k^2(XY^2/3-1/9)+1>0.
\]
Hence $\tau-\ell>k/3>s$.  For the second,
$v$ is a positive multiple of $4XY^2$, so
$v/M\ge XY/q^2=wY/q\ge wq^2>w$.
Substitution gives $Z=M_0\tau$ and $D_*=Mg$, establishing $A_0=A_1=0$.
This proves the reverse existence implication.  Composition with the
native syntactic calibration and negation gives the claimed equivalence
with $\operatorname{Con}(\mathrm{ZFC})$, by ordinary mathematics formalizable
in ZFC.

\subsection{Exact degree and finite verification}

The cancellations
\begin{align*}
 Z&=M_0\ell+Z_t,\\
 A_0&=2M_0\ell Z_t+Z_t^2-M_0^2(\ell k+1),\\
 (X+ac+H\gamma)^2-(1+\Delta c^2)
 &=2ac(X+H\gamma)+(X+H\gamma)^2-Hc^2-1,\\
 \widehat Y^{\,2}-\widehat V^{\,2}&=Lu(2\widehat V+Lu)
\end{align*}
give residual degrees $(13,13,13,14)$.  The degree-$14$ homogeneous part
of $A_s$ is
\[
 -9\beta^2i^2(\eta+\zeta)^4X^2Y^6.
\]
Only $A_s^2$ reaches degree $28$, so the leading form of $P_{13}$ is
\[
 \boxed{81\beta^4i^4(\eta+\zeta)^8X^4Y^{12}.}
\]
The coefficient of $i^4\eta^8X^4Y^{12}$ is $81\beta^4>0$ for the actual
fixed instance.  The exact degree is therefore unaffected by specialization
of the other three native ports.

The independent standard-library checker expands the four literal
residuals into respectively $1921$, $5565$, $26877$, and $81$ symbolic
monomials, compares separately specified formulas, checks their exact
witness degrees, and verifies the final sum of squares over independent
residual indeterminates.  It also checks all wires and liveness.  The
counts are $25$ additions, $16$ subtractions, and $57$ multiplications:
$98$ operations, or $91$ for the four-residual system.  The normal and
optimized Python runs both pass.  These finite checks certify the emitted
algebra and counts; they neither decide consistency nor replace the
unbounded mathematical argument and its named inherited dependencies.
